\section{Disk moments and deterministic coefficients}
\label{coef:section}

We now supply the deterministic numbers that select the gauge. First,
the physical-domain multipoint estimate proves that one centered disk
mass has positive finite moments. Then its conditional Green potential
recovers those moments through an explicit sequence of integrals.
The last subsection realizes the integrands by normal coordinates and
ordinary differential equation limits.

\subsection{The centered disk model}

For $0<\theta<\pi$, use the conformal identification
$J_\theta:\HH\to\DD$ below.  It takes $e^{i\theta}$ to the marked center
$0$.  Denote the ordinary law in these coordinates by $\PP_\theta$, and
its centered finite-level laws at $0$ by $\PP_\theta^*$.
Average the latter over the invariant angle law $\pi_a$, and write
$\PP_{\mathrm{mix}}^*$ for the resulting consistent stopped laws.
Define
\begin{equation}
 \begin{gathered}
 J_\theta(\zeta)=\frac{\zeta-e^{i\theta}}{\zeta-e^{-i\theta}},
 \qquad K=\{x\in\DD:|x|\le\tfrac14\},
 \qquad \PP_{\mathrm{mix}}^*=\int_0^\pi\PP_\theta^*\,\pi_a(\dd\theta),
 \\
 L_\ell=\mu^{\DD}_{\tau_\ell(0)}(K),
 \qquad L=\lim_{\ell\to\infty}L_\ell.
 \end{gathered}
 \label{coef:disk-model}
\end{equation}
The disk mass is the weighted pushforward of the single canonical
half-plane functional in Section~\ref{mass:section}.
Each $L_\ell$ is a raw finite-stopped-path quantity; its increasing limit
is defined on the consistent collection of such quantities.  In
particular \eqref{coef:disk-model} does not require a Brownian
continuation after the centered terminal time.

\begin{proposition}[Finite, positive disk moments]\label{coef:positive-moments}
For every integer $n\ge1$, the deterministic numbers
\begin{equation}
 m_n(\ell):=\EE_{\mathrm{mix}}^*L_\ell^n,
 \qquad
 m_n:=\EE_{\mathrm{mix}}^*L^n
       =\lim_{\ell\to\infty}m_n(\ell)
       \quad\text{satisfy}\quad 0<m_n<\infty.
 \label{coef:moments}
\end{equation}
Moreover $L>0$ almost surely, and, for each fixed angle,
\begin{equation}
 \sup_{\ell>0}\EE_\theta^*L_\ell^n
                  \le C_{n,a}\sin^{-p}\theta.
 \label{coef:angle-moments}
\end{equation}
\end{proposition}

\begin{proof}
We apply the multipoint estimate directly in the physical disk.
Its initial conformal radius is $R_0(x)=1-|x|^2$, independently of
$\theta$.  Choose a nonnegative continuous cutoff $f\ge\ind_K$ supported
in a fixed compact subdisk.  Apply the center-hit bound of Corollary~\ref{mp:level-moments} with
$m=n$, $x_0=0$, $e_0=e^{-2a\ell}$, and level weight
$A_j(f)=M_\infty^{(j),\DD}(f)$.  The proof of that corollary includes
$0$ as one extra fixed point in the multipoint estimate.  Its caps and
collision integrals can be chosen in the physical disk, where $K$,
$\supp f$, and the initial radii are independent of $\theta$.  Thus
\[
 \EE_\theta\!\left[
   \ind_{\{\tau_\ell(0)<\infty\}}
   \bigl(M_\infty^{(j),\DD}(f)\bigr)^n\right]
       \le C_{n,a}e^{-2a\alpha\ell},
\]
uniformly in $j,\ell,\theta$.

There is no angle-dependent extraction in taking this limit.
Every disk approximant is exactly the weighted pushforward of the
half-plane approximant.  Convergence of the fixed convex combinations
holds for every compact continuous half-plane test, and therefore for
the pullback of this $f$ under every fixed $J_\theta$.
Apply the displayed estimate before taking convex combinations;
convexity and Fatou's Lemma, with the center-hit event retained, yield
\begin{equation}
 \EE_\theta\bigl[
   \ind_{\{\tau_\ell(0)<\infty\}}\mu_\infty^{\DD}(K)^n\bigr]
       \le C_{n,a}e^{-2a\alpha\ell}.
 \label{coef:ordinary-hit-moment}
\end{equation}
Thus the uniform constant has not been obtained by transporting a
half-plane estimate through maps that degenerate near endpoint angles.

Since $R_0(0)=1$, the density at $\tau_\ell(0)$ is at most
$e^{2a\alpha\ell}\sin^{-p}\theta$ on the center-hit event.
Also $\mu_{\tau_\ell}^{\DD}(K)\le\mu_\infty^{\DD}(K)$ under the ordinary
law.  The finite-level change of law and
\eqref{coef:ordinary-hit-moment} prove \eqref{coef:angle-moments}.
The bound is integrable against $\pi_a$, because
\[
 \int_0^\pi\sin^{-p}\theta\,\pi_a(\dd\theta)
    =c_a\int_0^\pi\sin^{4a-p}\theta\,\dd\theta
    =c_a\int_0^\pi\sin\theta\,\dd\theta<\infty.
\]
Monotone convergence gives finite moments of $L$ of every integer order.

For positivity, choose a finite inner level with
$e^{-2a\ell}<1/4$.  Its attainment under the centered law is certain.
At that time, the distance from $0$ to the surviving domain boundary is
less than $1/4$.  The nearby boundary cannot be part of the original
unit circle, so the past trace has entered the interior of $K$.
Continuity gives a nontrivial time interval, contained in this finite
stopped piece, whose trace lies in the interior of $K$.
The positive-increment assertion of Proposition~\ref{mass:positive}
and the increment-support assertion of Proposition~\ref{mass:construction}
transfer to this stopped piece by its finite-level density.  They give
positive mass in $K$.  Hence $L>0$ almost surely, and $m_n>0$.
\end{proof}

\subsection{A Wiener-integral formula}

We use the discrete potential compensator appearing in
Lawler--Zhou \cite[Theorem~1 and its proof]{LawlerZhou2013}.
We prove its convergence at a random radius level and then combine it
with the finite-level Green weight to specify the centered moments.

Let $W$ be standard Wiener measure on
$C_0([0,\infty),\mathbb R)$.  For a fixed $\theta$ and driver $b$, solve
the ordinary Loewner equation with numerator $a$, using $J_\theta$ for
the disk identification, and form
\[
 \Psi_t(\theta,b)=\int_K G_t(x;\theta,b)\,\dd A(x),
 \qquad \tau_\ell=\tau_\ell(0),
 \qquad
 P_\ell(\theta,b)
 =\ind_{\{\tau_\ell<\infty\}}e^{2a\alpha\ell}
       \frac{\sin^p\theta_{\tau_\ell}(0)}{\sin^p\theta}.
\]
The last expression is zero on failure.  For paths $b,c$, let
$b\oplus_t c$ agree with $b$ through time $t$, and equal
$b(t)+c(\,\cdot-t)$ thereafter.  Put
\[
 \delta_h=2^{-h},\qquad N_h=h2^h,\qquad t_i=i\delta_h,
\]
and define the potential differences and their stopped cumulative sum by
\begin{align}
 q_{h,i}(\theta,b)
   &=\Psi_{t_i}(\theta,b)
       -\int\Psi_{t_{i+1}}(\theta,b\oplus_{t_i}c)\,\dd W(c),
       \label{coef:increments}\\
 D_h(\theta,b;\ell)
   &=\sum_{i=0}^{N_h-1}\ind_{\{t_{i+1}\le\tau_\ell\}}
                      q_{h,i}(\theta,b).
       \label{coef:discrete-mass}
\end{align}
Set undefined differences to zero.  For every fixed $\theta$, these
exceptions have Wiener measure zero by the compact first moment and
the conditional kernel.  All integrands are measurable in $\theta$ and
the paths.  Write $[x]_0^u=\min\{\max\{x,0\},u\}$.

The next estimate explains why a potential formula recovers mass at a
random endpoint. For fixed $\theta$, write $A_t=\mu_t^{\DD}(K)$ and
interpret $A_{\tau_\ell}$ as $A_\infty$ when $\tau_\ell=\infty$.

\begin{lemma}[Recovery at a random level]\label{coef:recovery}
For every fixed $\theta\in(0,\pi)$ and finite $\ell>0$,
\[
 D_h(\theta,\cdot;\ell)\longrightarrow A_{\tau_\ell}
 \quad\text{in }L^2(\PP_\theta).
\]
\end{lemma}

\begin{proof}
The process $A_t$ is continuous,
increasing, starts at zero, and has a finite terminal limit with
$\EE_\theta A_\infty^2<\infty$.  Equation~\eqref{mass:kernel} gives
\[
 \Psi_t=\EE_\theta[A_\infty-A_t\mid\mathcal F_t].
\]
Independent Brownian continuation identifies the Wiener integral in
\eqref{coef:increments} with its conditional expectation.
The tower property therefore yields
\[
 q_{h,i}=\EE_\theta[\Delta_{h,i}A\mid\mathcal F_{t_i}],
 \qquad \Delta_{h,i}A=A_{t_{i+1}}-A_{t_i}.
\]
Put $C_{h,j}=\sum_{i<j}q_{h,i}$.  The grid process
$A_{t_j}-C_{h,j}$ is a square-integrable discrete martingale starting
at zero. Its increments are centered conditional mass increments;
their variances are bounded by the second moments of $\Delta_{h,i}A$.
Orthogonality of its increments and the maximal inequality imply
\begin{equation}
 \EE_\theta\max_{j\le N_h}|C_{h,j}-A_{t_j}|^2
       \le4\EE_\theta\sum_{i<N_h}(\Delta_{h,i}A)^2
       \longrightarrow0.
 \label{coef:maximal-error}
\end{equation}
To justify the limit despite the growing time horizon, set
\[
 \omega_A(\delta)
   =\sup_{\substack{s,t\ge0\\|s-t|\le\delta}}|A_t-A_s|.
\]
Continuity and the finite terminal limit give
$\omega_A(\delta)\to0$ almost surely.  One can first use uniform
continuity on a fixed finite interval and then make the total variation
of the remaining tail small.  Moreover,
\[
 \sum_{i<N_h}(\Delta_{h,i}A)^2
       \le A_\infty\omega_A(\delta_h)\le A_\infty^2,
\]
so dominated convergence proves \eqref{coef:maximal-error}.

Let $j_h=\min\{N_h,\lfloor\tau_\ell/\delta_h\rfloor\}$, with
$j_h=N_h$ if $\tau_\ell=\infty$.  Then $D_h=C_{h,j_h}$.
On a finite hit $t_{j_h}\to\tau_\ell$; on failure $t_{j_h}=h\to\infty$.
Thus $A_{t_{j_h}}\to A_{\tau_\ell}$ in $L^2$, interpreting the latter as
$A_\infty$ on failure.  The maximum in \eqref{coef:maximal-error} controls
every selected index, so
\[
 D_h(\theta,b;\ell)\longrightarrow A_{\tau_\ell}
             \quad\text{in }L^2(\PP_\theta).
\]
No predictability or stopping-time property of $j_h$ is used.

\end{proof}

\begin{proposition}[Explicit coefficients]\label{coef:explicit}
The moments in \eqref{coef:moments} are given by
\begin{equation}
 \begin{aligned}
 m_n(\ell)
   &=\lim_{u\to\infty}\lim_{h\to\infty}
       \int_0^\pi\int P_\ell(\theta,b)
       \bigl([D_h(\theta,b;\ell)]_0^u\bigr)^n
       \,\dd W(b)\,\pi_a(\dd\theta),\\
 m_n&=\lim_{\ell\to\infty}m_n(\ell).
 \end{aligned}
 \label{coef:formula}
\end{equation}
All limit indices may be restricted to positive integers.  Thus these
coefficients depend only on $\kappa$ and the specified Gaussian, angle
and area integrals.
\end{proposition}

\begin{proof}
Fix $\theta,\ell,u$. The function $f_u(x)=([x]_0^u)^n$ is bounded
and continuous. Lemma~\ref{coef:recovery} gives
$f_u(D_h)\to f_u(A_{\tau_\ell})$ in probability. This convergence
passes through the fixed integrable density $P_\ell$: truncate this density at a constant $R$,
use bounded convergence in probability there, and then let $R\to\infty$.
Since $\int P_\ell\,\dd W=1$ for each fixed angle, the resulting
truncated Wiener integrals are bounded by $u^n$.  Dominated convergence
therefore permits the angle average, without a convergence estimate
uniform in $\theta$.  The finite-level change of law identifies the
limit as $\EE_{\mathrm{mix}}^*([L_\ell]_0^u)^n$.
Finally, monotone convergence in $u$ and then in $\ell$ proves
\eqref{coef:formula}.
\end{proof}

\subsection{A prescription using only normal coordinates and ODE limits}

The Wiener integrals in \eqref{coef:formula} are deterministic
integrals. The following construction specifies their integrands without
using a sample of the trace in Theorem~\ref{thm:main}.  They can be realized
as countable product standard-normal integrals: choose each unit
increment of $b$ as an independent standard normal, and recursively set
the midpoint of a dyadic interval of length $2^{-k}$ equal to the average
of its endpoint values plus $2^{-k/2-1}$ times a new independent standard
normal.  The polygonal interpolants converge uniformly on compacts almost
surely.  Indeed Gaussian tails and Borel--Cantelli bound the largest
level-$k$ normal on finitely many unit intervals by $O(\sqrt{k+1})$,
and $\sum_k2^{-k/2}\sqrt{k+1}<\infty$.  The limit has Gaussian covariance
$\min\{s,t\}$ and hence Wiener law.  Use the zero path on the null
exceptional set.

The Loewner quantities in the integrands do not require constructing the
trace.  For $\zeta=J_\theta^{-1}(x)$ and an integer $j\ge1$, use the
truncated vector field
\[
 F_j(t,z)=\frac{a\,\overline{z-b(t)}}
                    {\max\{|z-b(t)|^2,j^{-4}\}}.
\]
Starting at $\zeta$, perform the Euler iteration
$g_{k+1}=g_k+2^{-N}F_j(k2^{-N},g_k)$ and interpolate linearly.
For fixed $j$, its compact-uniform limit $g^{(j)}$ exists as $N\to\infty$:
the vector field is bounded, continuous in time, and globally Lipschitz
in its state variable.  A given finite time $t$ is surviving exactly when
some $j$ satisfies
\[
 \min_{0\le u\le t}\operatorname{Im}g_u^{(j)}>j^{-1}.
\]
For such a witness, $|g^{(j)}-b|^2>j^{-2}\ge j^{-4}$, so the cutoff is
inactive and uniqueness gives the genuine Loewner solution on $[0,t]$.
Conversely, a surviving solution on this compact interval has positive
minimum imaginary part, and a sufficiently large $j$ recovers it.

On survival use this solution to compute
\begin{align*}
 s_x(t)&=\int_0^t\frac{(\operatorname{Im}g_u)^2}{|g_u-b(u)|^4}\,\dd u,\\
 R_t(x)&=2\operatorname{Im}\zeta\,|J_\theta'(\zeta)|e^{-2as_x(t)},\\
 G_t(x)&=R_t(x)^{-\alpha}\sin^p\arg(g_t-b(t)),
\end{align*}
and put $G_t(x)=0$ off survival.  At the marked center the level time is
\[
 \tau_\ell=
 \inf\{q\in\mathbb Q_{>0}:q\text{ is surviving and }s_0(q)>\ell\},
 \qquad \inf\varnothing=\infty.
\]
Continuity and strict increase of the clock on its open lifetime identify
this infimum with attainment of level $\ell$ while alive.  These rational
tests, deterministic limits and integrals specify every integrand in
\eqref{coef:formula}.  This is an infinite-limit prescription; no finite computational error
bound is asserted. The coefficients retain neither an unspecified
normalization of trace mass nor a choice of convex extraction.
